\subsection{有限特征的性质}

定义2. 设E是一个集合。E的子集的一个集合S称为具有有限特征，如果关系X∈S等价于关系“X的每个有限子集都属于S”。

一个性质 \(P\{X\}\) 关于集合E的子集X，如果使 \(P\{X\}\) 成立的E的子集X的集合具有有限特征，则称该性质具有有限特征。

\minorheading{示例}

\begin{enumerate}
\def\labelenumi{(\arabic{enumi})}
\tightlist
\item
  一个有序集合E的全序子集的集合具有有限特征。事实上，E的一个子集X是全序的当且仅当X的每个由两个元素组成的子集都是全序的。
\end{enumerate}

*（2）一个模的所有自由子集的集合具有有限特征。一个域的扩张的所有代数自由子集的集合也是如此。

\begin{enumerate}
\def\labelenumi{(\arabic{enumi})}
\setcounter{enumi}{2}
\tightlist
\item
  一个模E的子模的集合不具有有限特征，因为E的一个子模的有限子集不一定是E的子模。*
\end{enumerate}

定理1. 集合E的子集的一个具有有限特征的集合S（按包含关系排序时）有一个极大元素。

由§2第4号的定理2，只需证明 \(\mathfrak{S}\) 是归纳的。为此我们将证明：如果 \(\mathfrak{G}\) 是 \(\mathfrak{S}\) 的任意一个按包含关系全序的子集，则 \(\mathfrak{G}\) 的诸集合的并集X属于 \(\mathfrak{S}\) (§2，第4号，定理2，推论2)。由于 \(\mathfrak{S}\) 具有有限特征，只需证明X的每个有限子集Y都属于 \(\mathfrak{S}\) 。现在，对于每个 \(y \in Y\) ，存在一个集合 \(Z_y \in \mathfrak{G}\) 使得 \(y \in Z_y\) 。由于集合族 \(Z_y\) ( \(y \in Y\) )是有限的且按包含关系全序排列，它有一个最大元素S（第4号，命题3的推论1）；换言之，存在一个集合 \(S \in \mathfrak{G}\) 使得 \(Y \subset S\) 。但由于 \(S \in \mathfrak{S}\) ，且由于Y是S的有限子集，我们有 \(Y \in \mathfrak{S}\) 因为 \(\mathfrak{S}\) 具有有限特征；这就完成了证明。

